\section{Definición y propiedades de la función de puntuación}

\subsection{Definición de la función de puntuación}

La función de puntuación (score function) tiene una definición precisa en estadística: es el gradiente del logaritmo de la función de densidad de probabilidad:

\begin{importantbox}{Definición de la función de puntuación}
Dada una función de densidad de probabilidad $p(x)$, su función de puntuación se define como:
$$s(x) = \nabla_x \log p(x)$$
donde $\nabla_x$ denota el gradiente respecto al punto de datos $x$ (y no respecto a los parámetros del modelo).
\end{importantbox}

Es importante notar que, en este contexto, la palabra ``puntuación'' no se refiere al concepto común de ``nota'' o puntaje, sino que es un término técnico con un significado específico en estadística. La función de puntuación es un campo vectorial: para cada punto $x$ en el espacio de datos, proporciona un vector de la misma dimensión que apunta hacia la dirección en la que la probabilidad logarítmica aumenta más rápidamente.

\begin{knowledgebox}{Comprensión intuitiva de la función de puntuación}
La función de puntuación $\nabla_x \log p(x)$ puede entenderse como una ``dirección guía'':
\begin{itemize}
    \item En regiones cercanas a alta densidad de probabilidad, el valor de la puntuación es pequeño (porque ya se está cerca de la ``cima'').
    \item En regiones de baja densidad de probabilidad, la puntuación apunta hacia la dirección de mayor densidad.
    \item Los puntos donde la puntuación es cero corresponden a los puntos extremos de la densidad de probabilidad.
\end{itemize}
En términos intuitivos, la función de puntuación nos indica ``hacia dónde debemos ir para llegar a una zona con mayor probabilidad''.
\end{knowledgebox}

\subsection{Función de puntuación de la distribución gaussiana}

Para establecer una intuición, calculemos la función de puntuación de una distribución gaussiana. Para la distribución condicional (es decir, la distribución del salto hacia adelante):

$$q(x_t \mid x_0) = \mathcal{N}(x_t;\; \sqrt{\bar{\alpha}_t}\, x_0,\; (1-\bar{\alpha}_t) I)$$

su función de puntuación es:

$$\nabla_{x_t} \log q(x_t \mid x_0) = -\frac{x_t - \sqrt{\bar{\alpha}_t}\, x_0}{1-\bar{\alpha}_t}$$

Este derivado es directo: primero se escribe la función de densidad de probabilidad de la distribución gaussiana, luego se toma el logaritmo y se retienen solo los términos relacionados con $x_t$ (es decir, el término cuadrático), y finalmente se calcula el gradiente respecto a $x_t$.

\begin{importantbox}{Relación entre predicción de ruido y función de puntuación}
Recordando que $x_t = \sqrt{\bar{\alpha}_t}\, x_0 + \sqrt{1-\bar{\alpha}_t}\, \epsilon_t$, se obtiene:
$$\epsilon_t = \frac{x_t - \sqrt{\bar{\alpha}_t}\, x_0}{\sqrt{1-\bar{\alpha}_t}}$$
Por lo tanto:
$$\nabla_{x_t} \log q(x_t \mid x_0) = -\frac{\epsilon_t}{\sqrt{1-\bar{\alpha}_t}}$$
Es decir, el $\epsilon_t$ predicho por el predictor de ruido en DDPM difiere de la función de puntuación condicional únicamente por un factor de escala. El significado físico de la predicción de ruido es aprender la función de puntuación (hasta un factor de escala).
\end{importantbox}

Este descubrimiento revela el significado profundo del predictor de ruido de DDPM: aunque superficialmente estamos ``prediciendo ruido'', en realidad estamos aprendiendo la función de puntuación de la distribución del salto hacia adelante.

\subsection{De la puntuación condicional a la puntuación marginal}

Hasta ahora hemos discutido la función de puntuación de la distribución \textbf{condicional} $q(x_t \mid x_0)$, ya que conocemos $x_0$. Sin embargo, en el proceso de generación no conocemos $x_0$; lo que realmente necesitamos es la función de puntuación de la distribución \textbf{marginal} $q(x_t)$:

$$\nabla_{x_t} \log q(x_t) = ?$$

donde $q(x_t) = \int q(x_t \mid x_0)\, q(x_0)\, dx_0$. Dado que $q(x_0)$ (la distribución real de los datos) suele ser desconocida y compleja, $q(x_t)$ generalmente no tiene una forma analítica cerrada.

No obstante, mediante deducciones matemáticas se puede demostrar:

$$\nabla_{x_t} \log q(x_t) = \mathbb{E}_{q(x_0 \mid x_t)}\left[\nabla_{x_t} \log q(x_t \mid x_0)\right]$$

\begin{importantbox}{La puntuación marginal es la esperanza a posteriori de la puntuación condicional}
La función de puntuación de la distribución marginal es igual al valor esperado de la función de puntuación condicional bajo la distribución a posteriori $q(x_0 \mid x_t)$. Esto significa que la red neuronal predictora de ruido $\epsilon_\theta(x_t, t)$ entrenada efectivamente está aprendiendo la función de puntuación obtenida promediando a posteriori todos los posibles $x_0$.
\end{importantbox}

\subsection{Resumen de la sección}

La función de puntuación se define como el gradiente de la densidad de probabilidad logarítmica y codifica información completa sobre la distribución de probabilidad. La función de puntuación de la distribución gaussiana tiene una solución cerrada concisa, mientras que el predictor de ruido en DDPM esencialmente aprende la función de puntuación de la distribución del salto hacia adelante (con un factor de escala). La puntuación marginal es igual a la esperanza de la puntuación condicional bajo la distribución a posteriori, lo que proporciona una explicación matemática rigurosa para entender el objetivo de aprendizaje de las redes neuronales en los modelos de difusión.
